\documentclass[11pt]{article}

\usepackage[margin=1in]{geometry}
\usepackage{amsmath,amssymb,amsthm}
\usepackage{booktabs}
\usepackage{enumitem}
\usepackage[colorlinks=true,linkcolor=blue,citecolor=blue,urlcolor=blue]{hyperref}

\newtheorem{theorem}{Theorem}[section]
\newtheorem{lemma}[theorem]{Lemma}
\newtheorem{proposition}[theorem]{Proposition}
\newtheorem{corollary}[theorem]{Corollary}
\theoremstyle{definition}
\newtheorem{definition}[theorem]{Definition}
\newtheorem{remark}[theorem]{Remark}
\newtheorem{question}[theorem]{Question}

% Honest evidence labels
\newcommand{\PROVED}{\textup{[\textsc{proved}]}}
\newcommand{\PROVEDVER}{\textup{[\textsc{proved + verified}]}}
\newcommand{\VERIFIED}{\textup{[\textsc{verified exhaustively}]}}
\newcommand{\COMPUTED}{\textup{[\textsc{computed exactly}]}}
\newcommand{\MEASURED}{\textup{[\textsc{measured}]}}

\newcommand{\Z}{\mathbb{Z}}
\newcommand{\N}{\mathbb{N}}
\newcommand{\vtwo}{v_2}
\newcommand{\vthree}{v_3}
\newcommand{\Info}{\mathrm{I}}
\newcommand{\Ent}{\mathrm{H}}

\title{The $3x+1$ map is a maximally forgetful machine:\\
exact fairness, zero storage, and the carry characterization}

\author{Martien de Jong\\[2pt]
\small with computational collaboration by Jengo (AI)}

\date{July 2026}

\begin{document}

\maketitle

\begin{abstract}
Conway proved in 1972 that generalized Collatz maps are computationally universal, and
Kurtz and Simon later located the generalized problem at $\Pi^0_2$-completeness. Both
constructions work by \emph{storing information}: register contents are encoded in the
exponents of primes and manipulated by the affine branches. We ask whether the specific
$3n+1$ instance shows any signature of such storage, and we report that every test we
can perform returns the opposite signature, in three layers. (1) \emph{Carry
characterization} (proved, elementary): in base $3$ the map $r \mapsto 3r+1$ is a pure
digit append, $\mathrm{trits}(3r+1) = \mathrm{trits}(r) \mathbin{+\!\!+} [1]$, and it is
the unique member of the $3n+d$ family with this property; every other offset carries or
borrows into the digits of $r$ with geometrically distributed depth
$P(\mathrm{depth}=j) = 2/3^j$. (2) \emph{Exact fairness} (proved via the Terras
correspondence, verified exhaustively): the density of odd $n$ whose first $r$
successive ladder refills all have depth at least $j$ is exactly $2^{-(j-1)r}$; and the
within-step joint law of (trailing ones $k$, post-burn halvings $w$) factorizes exactly,
$P(w \mid k) = 2^{-w}$ for every $k$. (3) \emph{Zero storage} (verified exhaustively):
the mutual information between successive refill depths along macro-orbits is
$0$ to within $3 \times 10^{-6}$ bits at distances $d = 1,2,3$, and the apparent
cross-base coupling $\Info(n \bmod 9;\, w_{\mathrm{next}})$ vanishes as the $2$-adic
modulus grows. At the operator level, the linearized Krasikov--Lagarias edge system has
spectrum $\{1\} \cup \{\,|\lambda| = 1/4\,\}$ exactly (flat damping), and the identity
$2^\alpha = 3$ pins its average row mass to exactly $1$ (conservative line). The
density-level stochastic model of the $3n+1$ map is therefore \emph{exactly} independent
fair geometric dice with no measurable correction --- the opposite of a hidden computer.
We discuss what this anti-universality evidence does and does not say about the
decidability of the specific instance, and we prove that density-level arguments are
map-blind: the growth exponent of the Krasikov--Lagarias method is identical for $3n-1$
and $3n+1$ by mirror symmetry, so any proof of the conjecture must use the carry
characterization --- the only lever that distinguishes the two maps.
\end{abstract}

\section{Motivation: where could a computer hide?}

The Collatz map on odd integers, in Syracuse form
\[
T(n) \;=\; \frac{3n+1}{2^{\vtwo(3n+1)}},
\]
is conjectured to send every positive odd integer to $1$. Two classical results frame
the pessimistic view of this problem as a question about \emph{computation} rather than
arithmetic.

First, Conway \cite{Conway1972} showed that the natural generalization --- piecewise
affine maps $g(n) = a_i n + b_i$ for $n \equiv i \pmod{p}$ --- is computationally
universal: one can encode any register machine in such a map, and the question ``does
the orbit of $n$ reach $1$?'' is undecidable for the class. Second, Kurtz and Simon
\cite{KurtzSimon2007} sharpened this to $\Pi^0_2$-completeness: the statement ``every
orbit of $g$ eventually reaches $1$'' is, for generalized Collatz maps, as hard as any
$\forall\exists$ statement of arithmetic can be.

Both constructions have a common engine: \emph{information storage}. The state of a
Minsky machine's registers is encoded in the exponents of small primes dividing $n$; the
affine branches increment, decrement, and test those exponents. A generalized Collatz
map is undecidable exactly because, and exactly insofar as, it can \emph{remember}.

This suggests a concrete empirical program for the specific instance $3n+1$. If the
map were secretly a computer --- if the hardness of the conjecture were of Conway's kind
--- we would expect measurable traces: correlations between successive branch decisions,
residue channels that carry information forward, deviations from the naive stochastic
model that a memory mechanism would necessarily produce. Conversely, if every such test
returns exact independence, exact fairness, and zero channel capacity, then the specific
instance carries the \emph{opposite} signature: not a register machine but an odometer,
a machine that forgets everything it does.

This paper collects six results from our research program (de Jong--Jengo, 2026,
research notes \cite{Note2026}) that carry out this program. Throughout we use honest
evidence labels: \PROVED{} means an unconditional proof is given or referenced;
\VERIFIED{} means an exhaustive computation over a stated finite range;
\COMPUTED{} means an exact finite computation (e.g.\ of a spectrum); \MEASURED{} means
a sampled or fitted quantity. The distinction matters and we never blur it.

\subsection*{Summary of results}

\begin{enumerate}[label=(\roman*)]
\item \textbf{Carry characterization} (Theorem~\ref{thm:carry}, \PROVEDVER):
$3r+1$ is the unique base-$3$-local (carry-free, append-only) member of the $3n+d$
family; all other offsets propagate into the digits of $r$ with geometric depth
$P(\mathrm{depth}=j) = 2/3^j$.
\item \textbf{Hop-tax / exact fairness} (Theorem~\ref{thm:hoptax}, \PROVEDVER): the
density of odd $n$ whose first $r$ ladder refills are all $\geq j$ is exactly
$2^{-(j-1)r}$.
\item \textbf{Zero storage across steps} (Theorem~\ref{thm:zerostorage}, \VERIFIED):
$\Info(k_0; k_d) = 0.000000$ bits for $d = 1,2,3$.
\item \textbf{Exact factorization within steps} (Theorem~\ref{thm:factorization},
\VERIFIED): $\Info(k;w) = 0$ and $P(w \mid k) = 2^{-w}$ for every $k$; apparent
cross-base coupling is finite-modulus leakage.
\item \textbf{Flat damping} (Theorem~\ref{thm:spectrum}, \COMPUTED{}, and
Lemma~\ref{lem:mass}, \PROVED): the linearized Krasikov--Lagarias edge operator has
spectrum $\{1\} \cup \{|\lambda| = 1/4\}$, and mass conservation
$3W_0 + W_2 + W_8 = 3$ holds exactly because $2^\alpha = 3$.
\item \textbf{Map-blindness of density arguments} (Proposition~\ref{prop:mirror},
\PROVED): $\lambda^*(3n-1) = \lambda^*(3n+1)$ by mirror symmetry; density methods
cannot separate the convergent map from its cycle-rich sibling.
\end{enumerate}

\section{Setup and notation}

For odd $n \geq 1$ write
\[
n + 1 = a \cdot 2^k, \qquad a \text{ odd},
\]
so that $k = k(n)$ is the number of trailing $1$s in the binary expansion of $n$
(the \emph{refill depth} or \emph{sequence index}) and $a$ is the \emph{family}. The
\emph{segment} is the run of $k$ exact odd steps $n \mapsto (3n+1)/2$, which ends at
$a \cdot 3^k - 1$ (an elementary computation; see \cite{Note2026}, Theorem~1). The
\emph{macro-step} appends the full cascade of
\[
w \;=\; \vtwo\!\left(a \cdot 3^k - 1\right)
\]
halvings, landing on the next odd number, whose own trailing-ones count $k'$ is the
\emph{next refill}. An orbit is thus coded by its \emph{index stream}
$(k_0, w_0), (k_1, w_1), \dots$, and the naive stochastic model --- going back to
Terras \cite{Terras1976} and surveyed by Lagarias \cite{Lagarias1985, Lagarias2010} ---
treats $k$ and $w$ as independent geometric$(1/2)$ variables. The multiplicative drift
per macro-step under this model is
$\mathbb{E}[\log_2(3^k/2^{k+w})] = 2\log_2 3 - 4 < 0$, which is the heuristic reason
orbits fall.

The empirical content of this paper is that, at the level of natural density, the naive
model is not an approximation. It is \emph{exact}, with no measurable correction of any
kind, at every level we can test.

A word on the mechanism of exactness. Terras's correspondence \cite{Terras1976} shows
that the first $t$ parity decisions of an orbit are determined by, and determine,
$n \bmod 2^t$: each depth-$t$ index pattern occupies \emph{exactly one} residue class
modulo $2^{(\text{bits used})}$. Pattern densities are therefore exact dyadic rationals,
not asymptotic estimates. This is what allows several statements below to be proved,
not merely measured; the exhaustive counts then serve as verification of the proofs and
of our bookkeeping.

\section{The carry characterization of $3n+1$}

The first result explains \emph{why} the $3n+1$ map, among all its siblings, is the one
with exactly solvable statistics. Write $\mathrm{trits}(r)$ for the base-$3$ digit
string of $r \geq 1$, and $\mathbin{+\!\!+}$ for string concatenation.

\begin{theorem}[Carry characterization; \PROVEDVER]\label{thm:carry}
\leavevmode
\begin{enumerate}[label=(\alph*)]
\item For every $r \geq 1$,
\[
\mathrm{trits}(3r+1) \;=\; \mathrm{trits}(r) \mathbin{+\!\!+} [1]:
\]
the map $r \mapsto 3r+1$ is a pure append in base $3$, with zero propagation into the
digits of $r$.
\item Among the maps $r \mapsto 3r + d$ with $d$ odd and $3 \nmid d$, the offset
$d = 1$ is the unique one with this property. For every other such $d$ the map carries
or borrows into the digits of $r$.
\item \textup{(Depth law, proved for $d = -1$.)} For $d = -1$,
\[
\mathrm{trits}(3r-1) \;=\; \mathrm{trits}(r-1) \mathbin{+\!\!+} [2],
\]
and the borrow propagates into $r$ to depth $j$ (i.e.\ alters exactly the last $j$
trits of $r$) with natural density
\[
P(\mathrm{depth} = j) \;=\; \frac{2}{3^j}, \qquad j \geq 1.
\]
\end{enumerate}
\end{theorem}

\begin{proof}
(a) Multiplying by $3$ appends the digit $0$ in base $3$; adding $1$ converts that
trailing $0$ to $1$ with no carry, since $1 < 3$.

(b) Write $d = 3q + s$ with $s \in \{1, 2\}$ (possible since $3 \nmid d$). Then
$3r + d = 3(r + q) + s$, i.e.\ the map appends the digit $s$ to the digits of $r + q$.
This is an append onto the digits of $r$ itself if and only if $q = 0$, i.e.\
$d \in \{1, 2\}$; among odd offsets, only $d = 1$. For every other $d$ the addition of
$q \neq 0$ to $r$ propagates a carry (if $q > 0$) or a borrow (if $q < 0$) into the
digit string of $r$.

(c) The identity is (b) with $q = -1$, $s = 2$: $3r - 1 = 3(r-1) + 2$. Passing from
$r$ to $r - 1$ borrows through exactly the trailing run of $0$s in
$\mathrm{trits}(r)$: if the last $j - 1$ trits of $r$ are $0$ and the $j$-th from the
end is nonzero, the borrow alters exactly the last $j$ trits. For $r$ uniform in a
range $[1, 3^m)$ and $j \le m$, the last $j-1$ trits are $0$ and the $j$-th is nonzero
with probability $(1/3)^{j-1} \cdot (2/3) = 2/3^j$; the density statement follows by
letting the range grow.
\end{proof}

\emph{Verification.} Part (c)'s depth law was confirmed to $4+$ decimals by direct
count (\cite{Note2026}, Theorem~25 and \texttt{scripts/}). For general $|q| > 1$ the
carry is again confined to a trailing-digit run and is geometrically distributed with
the same $3^{-j}$ scale; we have proved the exact constants only for $|q| = 1$ and
measured the rest, so the general-$d$ depth constants carry the label \MEASURED.

\begin{remark}[The fingerprint]
The carry structure is not a curiosity; it is visible in the dynamics. The disorder
injected at ternary digit $j$ by the $3n-1$ borrow (probability $\sim 3^{-j}$) shows up
as a measured difference in the cascade heterogeneity profiles of the two maps:
$0.178$ versus $0.059$ at digit $1$ (\MEASURED, \cite{Note2026}, Theorem~25). The
convergent map is the maximally \emph{local} one: its ternary clockwork is
memoryless-append, which is why its stationary ``roulette'' measure is exactly solvable
and its refills exactly fair (Section~\ref{sec:fair}). Its siblings pay for their
offset with a geometric-depth scribble into their own digits --- which is, in embryo,
the write head of Conway's storage mechanism.
\end{remark}

\section{Exact fairness: the hop tax}\label{sec:fair}

After a full burn (segment of $k$ odd steps plus its cascade of $w$ halvings), the orbit
lands on a new odd number with a new refill depth $k'$. Call an odd number \emph{rich at
level $j$} if $k \geq j$. The naive model predicts
$P(k \geq j) = 2^{-(j-1)}$ and predicts that successive refills are independent. Both
statements are exactly true at density level.

\begin{theorem}[Hop tax; exact refill fairness; \PROVEDVER]\label{thm:hoptax}
Fix $j \geq 2$ and $r \geq 1$. The set of odd $n$ whose first $r$ successive ladder
refills (the trailing-ones depths after each of the first $r$ full burns) are all
$\geq j$ is a finite union of residue classes, and its natural density is exactly
\[
\left(2^{-(j-1)}\right)^{r} \;=\; 2^{-(j-1)r}.
\]
\end{theorem}

\begin{proof}[Proof sketch]
By the Terras correspondence, the event ``the first $r$ refills are all $\geq j$''
is determined by finitely many binary digits of $n$ and is a union of residue classes
modulo $2^t$ for a suitable $t$; its density is therefore an exact dyadic rational. The
count of classes follows from the fact that each macro-step's refill condition
$k_i \geq j$ constrains exactly $j - 1$ fresh binary digits (the digits revealed by the
current burn), independently of the digits already consumed: conditioned on the entire
past pattern, the residue-class measure of $\{k_i \geq j\}$ is exactly $2^{-(j-1)}$.
Multiplying over $i = 0, \dots, r-1$ gives the claim.
\end{proof}

\emph{Verification.} Exhaustive count over all odd $n < 2^{22}$, for all
$j \in \{2, 3, 4\}$ and $r \in \{1, 2, 3\}$: each of the nine measured densities agrees
with $2^{-(j-1)r}$ to at least four decimal places
(\cite{Note2026}, Theorem~20; \texttt{scripts/59\_hop\_tax.py}). The exact target
values are:

\begin{center}
\begin{tabular}{lccc}
\toprule
density $= 2^{-(j-1)r}$ & $r = 1$ & $r = 2$ & $r = 3$ \\
\midrule
$j = 2$ & $0.5$    & $0.25$      & $0.125$ \\
$j = 3$ & $0.25$   & $0.0625$    & $0.015625$ \\
$j = 4$ & $0.125$  & $0.015625$  & $0.001953125$ \\
\bottomrule
\end{tabular}
\end{center}

\begin{corollary}[Caste memorylessness]
$P(\mathrm{rich} \to \mathrm{rich})$, i.e.\ the conditional density of $k_{i+1} \geq 2$
given $k_i \geq 2$, equals $1/2$ exactly --- not approximately. Being rich buys no
persistence: the hop tax is collected in full at every step.
\end{corollary}

\begin{remark}[What remains open]
Theorem~\ref{thm:hoptax} is a statement about natural density --- about \emph{almost
every} starting point in the sense of residue-class measure. A divergent orbit would
have to beat this exactly fair coin \emph{forever, pointwise}. The gap between the
density statement (proved) and the pointwise statement (open) is precisely the remaining
content of the Collatz conjecture; nothing in this paper closes it, and we flag every
place where the two could be confused.
\end{remark}

\section{Zero storage}

A computer needs a channel from its past to its future. We now measure the capacity of
every channel we can define at the macro-step level, and find them all exactly closed.

\begin{theorem}[Zero storage across macro-steps; \VERIFIED]\label{thm:zerostorage}
Along macro-orbits started from all odd $n < 2^{21}$ (exhaustive), the mutual
information between the initial refill depth and the refill depth $d$ macro-steps later
satisfies
\[
\Info(k_0;\, k_d) \;=\; 0.000000 \text{ bits} \quad (\text{within } 3 \times 10^{-6}
\text{ bits}) \qquad \text{for } d = 1, 2, 3,
\]
against an unconditional refill entropy $\Ent(k) = 1.9375$ bits. The macro-automaton
has zero measurable channel capacity at density level: it stores nothing.
\end{theorem}

\begin{theorem}[Exact factorization of the macro-step law; \VERIFIED]\label{thm:factorization}
Exhaustively over all odd $n < 2^{22}$:
\begin{enumerate}[label=(\alph*)]
\item the within-step joint law of $(k, w)$ factorizes exactly:
$\Info(k; w) = 0.000000$ bits, and
\[
P(w \mid k) \;=\; 2^{-w} \quad \text{to four decimals, for every value of } k;
\]
\item \textup{(cross-base blindness)} the apparent coupling between the ternary state
and the next cascade, $\Info(n_T \bmod 9;\, w_{\mathrm{next}})$, tends to $0$ as the
$2$-adic truncation modulus grows: raising the modulus from $2^{12}$ to $2^{16}$
drops the measured mutual information by a factor of about $70$. The residual coupling
at any finite modulus is finite-modulus leakage, not structure.
\end{enumerate}
Consequently the density-level stochastic model of the $3n+1$ map is exactly
\[
k \sim \mathrm{geom}(1/2), \qquad w \sim \mathrm{geom}(1/2), \qquad
\text{all independent, memoryless, cross-base blind,}
\]
with no measurable correction at any level tested.
\end{theorem}

\emph{Provenance.} Theorem~\ref{thm:zerostorage} is \cite{Note2026}, Theorem~22;
Theorem~\ref{thm:factorization} is \cite{Note2026}, Theorem~27. Part (a) of
Theorem~\ref{thm:factorization} admits the same Terras-style proof mechanism as
Theorem~\ref{thm:hoptax} (the digits determining $w$ are disjoint from and unbiased by
the digits determining $k$), and we expect it to be upgradable to \PROVED{} by the
residue-class argument; we conservatively label it as verified, since the note records
it as an exhaustive computation. Part (b) is intrinsically a limit statement measured
along a modulus sequence and carries the label \MEASURED{} for the rate, \VERIFIED{}
for the counts at each fixed modulus.

\begin{remark}[Why this is the anti-Conway signature]
In Conway's construction \cite{Conway1972} and in the Kurtz--Simon completeness proof
\cite{KurtzSimon2007}, the orbit's future depends on unboundedly old information:
register contents ride along inside the iterate, and the map reads them. Any such
mechanism, translated to the $3n+1$ coordinate system, would have to show up as one of:
(i) correlation between refill depths at distance $d$ (Theorem~\ref{thm:zerostorage}
says: none, to $3 \times 10^{-6}$ bits); (ii) coupling between the burn size and the
cascade that follows it (Theorem~\ref{thm:factorization}(a) says: exactly none); or
(iii) a channel between the ternary representation --- where the multiplication by $3$
lives --- and the binary decisions --- where the halvings live
(Theorem~\ref{thm:factorization}(b) says: leakage only, vanishing with the modulus).
Every register we can look for is empty.
\end{remark}

\section{Flat damping and the conservative line}

The results so far concern the forward statistics of orbits. The Krasikov--Lagarias
(K--L) difference-inequality method \cite{KrasikovLagarias2003} probes the same machine
from the backward direction: it propagates counting inequalities over residue classes
modulo $3^j$ and extracts a growth exponent $\gamma$ with
$\pi_1(x) \geq x^{\gamma}$ for the count of integers below $x$ that reach $1$. The
linearization of this system at its spectral edge gives an independent, operator-level
test of forgetfulness.

\begin{theorem}[Flat spectrum of the linearized edge operator; \COMPUTED]\label{thm:spectrum}
Linearize the K--L system at the edge $\lambda = 2$ (replacing the min-operation by the
triple mean) on the classes modulo $3^j$. For $j = 4, 5, 6$ (that is, on $27$, $81$,
and $243$ classes), the spectrum of the resulting operator is exactly
\[
\{1\} \;\cup\; \{\,\lambda : |\lambda| = \tfrac14\,\},
\]
i.e.\ a single leading eigenvalue $1$ (the exactly solvable stationary ``roulette''
direction) and \emph{all} other eigenvalues of modulus exactly
$1/4 = \lambda^{-2}\big|_{\lambda = 2}$.
\end{theorem}

\emph{Provenance.} \cite{Note2026}, Theorem~21; \texttt{scripts/60\_contraction\_spectrum.py}.
The computation is exact linear algebra on finite matrices; the label \COMPUTED{} means
it is a finite exact computation at three depths, not an asymptotic theorem for all $j$.

The flat spectrum says: apart from the single conserved direction, \emph{every} mode of
deviation from the stationary measure is damped at the uniform rate $1/4$ per level.
There are no slow modes, no near-critical directions in which information about the past
could survive --- linear damping is overwhelming and structureless. (The slow
$\sim 0.93$/digit decay of the measured tempering amplitude is therefore \emph{not} a
linear mode; it is sustained by per-level injection through the min-nonlinearity, which
is a separate, localized question \cite{Note2026}.)

\begin{lemma}[Mass conservation at the edge; \PROVED]\label{lem:mass}
At the edge $\lambda = 2$, with $\alpha = \log_2 3$, the K--L row masses satisfy
exactly
\[
W_0 + W_2 = 1 \ (\text{type } 2 \bmod 9), \qquad
W_0 = \tfrac14 \ (\text{type } 5), \qquad
W_0 + W_8 = \tfrac74 \ (\text{type } 8),
\]
whence
\[
3W_0 + W_2 + W_8 \;=\; \tfrac34 + 3 \cdot 2^{\alpha - 2} \;=\; \tfrac34 + \tfrac94
\;=\; 3,
\qquad \text{since } 2^{\alpha} = 3:
\]
the \emph{average} row mass is exactly $1$. Consequently, for the lattice recurrence
$\mathrm{CV}_P = a\,\mathrm{CV}_{P-1} + c\,\mathrm{CV}_{P+1}$ governing the
heterogeneity cascade of the K--L eigenvector, the coefficients obey
\[
a + c \;\leq\; 1 \qquad \text{(subcriticality; measured value } 0.9955 \text{ at }
k = 20\text{)}.
\]
\end{lemma}

\emph{Provenance.} \cite{Note2026}, Lemma~24. The row-mass identities and the average
being exactly $1$ are elementary consequences of the edge equations; the passage to
$a + c \leq 1$ uses the offset algebra of the exact difference identity
(\cite{Note2026}, Theorem~16) plus the triangle inequality. The numerical value
$0.9955$ and the interpretation of the slack as a measured incoherence factor are
\MEASURED.

\begin{remark}[The defining equation pins the system to the conservative line]
The identity $2^{\alpha} = 3$ --- which is nothing but the definition of the problem,
the exchange rate between the doubling and tripling that constitute the map --- is
precisely what forces the average row mass to equal $1$. The system sits exactly
\emph{on} the conservative line $a + c = 1$: the heterogeneity cascade can never grow
exponentially, and the only open question at this level is the drift direction along
the line ($\theta = a/c < 1$), not the possibility of explosion. A hidden computer
would need super-critical mass somewhere to amplify and preserve a signal; the defining
equation of the map forbids it on average, and the flat $1/4$ spectrum
(Theorem~\ref{thm:spectrum}) crushes whatever remains.
\end{remark}

\section{Synthesis: the opposite of a hidden computer}

Assembling the layers:

\begin{itemize}
\item \textbf{Digit level} (Theorem~\ref{thm:carry}, proved): the map is a carry-free
append in base $3$ --- the unique such member of its family. There is no write-back
into the multiplicand's digits; the ternary tape only ever grows by one clean symbol.
\item \textbf{Step level} (Theorems~\ref{thm:hoptax} and \ref{thm:factorization}):
refills are exactly fair geometric draws, exactly independent of each other and of the
cascades, with the cross-base channel closed in the limit.
\item \textbf{Operator level} (Theorem~\ref{thm:spectrum}, Lemma~\ref{lem:mass}):
uniform flat damping at rate $1/4$ off a single conserved direction, with average mass
pinned to the conservative line by $2^{\alpha} = 3$ itself.
\end{itemize}

At every level, the density-level stochastic model is \emph{exactly} independent fair
geometric dice, with no measurable correction. This is the maximally forgetful point of
the $3n+d$ family, and Theorem~\ref{thm:carry} identifies the mechanism: carry-freeness.
The other members of the family scribble geometric-depth corrections into their own
digits --- a rudimentary memory --- and Conway's universal machines are the fully
engineered endpoint of that same mechanism. The $3n+1$ instance sits at the opposite
endpoint: the unique offset with no write head at all.

\subsection{What this says about decidability --- and what it does not}

We state the epistemic status plainly.

\begin{enumerate}[label=(\arabic*)]
\item \textbf{This is evidence, not proof.} Undecidability of the generalized problem
\cite{Conway1972, KurtzSimon2007} does not imply anything about the specific instance,
and conversely no finite battery of independence tests can \emph{prove} that the
instance is decidable or that the conjecture is true. A pointwise conspiracy ---
a single measure-zero thread of integers beating the fair coin forever --- is exactly
what the density results do not exclude, and it is the entire remaining content of the
problem (\cite{Note2026}, Theorems~20, 27).
\item \textbf{The evidence is nonetheless directional.} Every known route to
undecidability in this family runs through storage. The measured channel capacity of
the $3n+1$ macro-automaton is zero to $3 \times 10^{-6}$ bits, its transition
statistics are exactly the memoryless ideal, and its linearized backward operator has
no slow modes. If the specific instance were hard in Conway's sense, the hardness would
have to live entirely in a density-zero, information-theoretically invisible thread ---
a logical possibility, but one with no positive evidence and with every measurable
proxy returning the opposite sign. We call this the \emph{anti-universality} reading:
the conjecture's remaining content is that a countable integer thread cannot exploit a
channel that provably has no capacity at density level.
\item \textbf{The pointwise/density gap is the problem.} We emphasize once more:
Theorems~\ref{thm:hoptax}--\ref{thm:factorization} constrain almost every orbit, in
residue-class measure. The conjecture is a statement about \emph{every} orbit. Nothing
above closes that gap; what the results do is \emph{localize} it, by closing every
hiding place for density-level structure. The content of the conjecture is now
irreducibly pointwise.
\end{enumerate}

\subsection{Density arguments are provably map-blind}

Finally, the carry characterization has a sharp methodological consequence: it is the
\emph{only} lever a proof can use, because every density-level functional is blind to
the difference between $3n+1$ and its cycle-rich mirror $3n-1$.

\begin{proposition}[Mirror symmetry; \PROVED]\label{prop:mirror}
Negation conjugates the two maps: writing $T_{\pm}(n) = (3n \pm 1)/2^{\vtwo(3n \pm 1)}$
on odd integers, we have $T_{-}(n) = -\,T_{+}(-n)$. Hence the $3n-1$ dynamics on the
positive odd integers is the $3n+1$ dynamics on the negative odd integers, up to sign.
Since the Terras residue-class structure and the Krasikov--Lagarias class system are
invariant under this conjugation (it permutes residue classes modulo $2^t$ and modulo
$3^j$), every density-level quantity coincides for the two maps; in particular the K--L
growth exponent satisfies
\[
\lambda^{*}(3n-1) \;=\; \lambda^{*}(3n+1).
\]
\end{proposition}

\begin{proof}[Proof sketch]
$3(-n) + 1 = -(3n - 1)$ and $\vtwo(-m) = \vtwo(m)$, giving the conjugation identity.
The map $n \mapsto -n$ is a bijection on odd residues modulo $2^t$ and modulo $3^j$
that intertwines the branch structure of the two inverse iterations; densities of
residue-class unions are preserved. Any functional built from residue-class densities
--- pattern densities, predecessor-count inequalities, the K--L eigenvalue --- is
therefore identical for the two maps.
\end{proof}

But the two maps behave completely differently: $3n+1$ is conjectured (and
overwhelmingly verified) to funnel everything to $1$, while $3n-1$ on the positive
integers has at least three cycles ($\{1, 2\}$, the cycle through $5$, and the cycle
through $17$ --- equivalently the $-1$, $-5$, $-17$ cycles of $3n+1$ on the negatives;
cf.\ the census in \cite{Note2026}, Theorem~17). Proposition~\ref{prop:mirror} says no
density argument can see this difference. Where the maps \emph{do} differ is exactly
Theorem~\ref{thm:carry}: $+1$ appends, $-1$ borrows, with the measured disorder
fingerprint $0.178$ vs $0.059$ at ternary digit $1$. Therefore:

\begin{quote}
\emph{Any proof of the Collatz conjecture must use the carry-freeness of $3n+1$ ---
the sign of the offset in base 3 --- because it is the only property, at any level we
have found, that distinguishes the convergent map from its non-convergent mirror. A
method that would work verbatim for $3n-1$ is provably insufficient.}
\end{quote}

This is a filter of the same kind as the classical observation that a proof must
somehow use the archimedean size of $n$ and not just congruence information (cf.\ the
residue-blind impossibility result, \cite{Note2026}, Theorem~11): it does not tell us
how to prove the conjecture, but it tells us cheaply, in advance, which proposals cannot
work.

\subsection*{Open problems}

\begin{enumerate}[label=(\Alph*)]
\item Upgrade Theorem~\ref{thm:factorization}(a) from exhaustive verification to a
residue-class proof, and prove the cross-base limit in
Theorem~\ref{thm:factorization}(b) (leakage $\to 0$) at an explicit rate.
\item Prove the flat-spectrum statement of Theorem~\ref{thm:spectrum} for all depths
$j$, not just $j = 4, 5, 6$.
\item Quantify anti-universality: define a class of ``storing'' $3$-adic extensions of
the $3n+d$ family in which Conway's construction can be graded, and prove that
carry-free members have zero storage capacity as a theorem about the class, with
$3n+1$ as the unique carry-free point.
\item The pointwise problem itself, now stripped of every density-level disguise:
show that no single positive integer thread realizes an infinite net-climbing index
stream against an exactly fair, exactly memoryless background.
\end{enumerate}

\subsection*{Reproducibility}

All exhaustive counts and spectra referenced above are produced by short standalone
scripts in the research repository (\texttt{scripts/59\_hop\_tax.py},
\texttt{scripts/60\_contraction\_spectrum.py}, and the inline computations recorded in
\cite{Note2026}); ranges are stated in each theorem. No result in this paper depends on
sampling: every ``verified'' claim is an exhaustive computation over the stated range,
and every ``proved'' claim has an elementary proof included or referenced.

\begin{thebibliography}{9}

\bibitem{Conway1972}
J.~H. Conway,
\emph{Unpredictable iterations},
in: Proceedings of the 1972 Number Theory Conference,
University of Colorado, Boulder, CO, 1972, pp.~49--52.

\bibitem{KurtzSimon2007}
S.~A. Kurtz and J.~Simon,
\emph{The undecidability of the generalized Collatz problem},
in: Theory and Applications of Models of Computation (TAMC 2007),
Lecture Notes in Computer Science 4484, Springer, 2007, pp.~542--553.

\bibitem{Lagarias1985}
J.~C. Lagarias,
\emph{The $3x+1$ problem and its generalizations},
Amer. Math. Monthly 92 (1985), 3--23.

\bibitem{Lagarias2010}
J.~C. Lagarias (ed.),
\emph{The Ultimate Challenge: The $3x+1$ Problem},
American Mathematical Society, Providence, RI, 2010.

\bibitem{KrasikovLagarias2003}
I.~Krasikov and J.~C. Lagarias,
\emph{Bounds for the $3x+1$ problem using difference inequalities},
Acta Arith. 109 (2003), 237--258.

\bibitem{Terras1976}
R.~Terras,
\emph{A stopping time problem on the positive integers},
Acta Arith. 30 (1976), 241--252.

\bibitem{Note2026}
M.~de Jong and Jengo,
\emph{Six results on the Collatz map in family/pair coordinates},
research notes (2026), repository \texttt{collatz/research}, file \texttt{NOTE.md}
(Theorems 20, 21, 22, 25, 27; Lemma 24; and scripts cited therein).

\end{thebibliography}

\end{document}
